\documentclass[11pt]{article}

\usepackage[margin=1.1in]{geometry}
\usepackage[T1]{fontenc}
\usepackage{lmodern}
\usepackage{microtype}
\usepackage{booktabs}
\usepackage{fancyvrb}
\usepackage{upquote}
\usepackage[hidelinks,bookmarks=true,bookmarksopen=true,bookmarksdepth=3]{hyperref}

\fvset{fontsize=\small,frame=none,xleftmargin=1.5em}

\title{\normalsize\textsc{Whitepaper}\\[0.6em]
\LARGE surf: Context Traversal without Ingestion\thanks{Marc McCahill is
credited with coining the phrase ``surfing the Internet.'' This tool lets agents
surf the seas of context.}}
\author{Alexander R. Saint Croix\\\texttt{alex@saintx.us}}
\date{September 2026}

\begin{document}

\maketitle

\pdfbookmark[1]{abstract}{abstract}
\begin{abstract}
surf is a command-line tool that returns one section of a markdown, TeX, or PDF
document by heading, and returns the document's heading tree when no heading is
given. I built it so that agent skills can be written as thin indexes: a
\texttt{SKILL.md} that lists the intents an agent might arrive with and, for each
one, a single command that pulls the exact section of reference material that
intent needs. The reference corpus can then grow without the index changing, an
agent loads only what its task requires, and the relationship between index and
corpus can be checked by machine. The surf skill ships this paper as its TeX and
PDF example, so every command it describes can be run against it.
\end{abstract}

\section{Background}

Three approaches to giving a language model context have a defect in common.
Each tends to hand the model more than the task needs and leave it to sort out
which part applies.

Retrieval-augmented generation (RAG) selects chunks of a corpus by their
resemblance to a query and places them in the prompt. Resemblance to a query and
relevance to the agent's task are different measures, and RAG can compute only
the first. The model receives passages that look like the question and then
spends attention deciding which of them bear on the work.

Long instruction files are the second approach. A \texttt{CLAUDE.md} or
\texttt{AGENTS.md} loads on every turn, and a file that has grown to cover every
situation a project might present is, on any given turn, mostly about other
situations. The model attends over all of it to find the part that applies.

Skills are the third, and they improve on this by adding agency. A skill loads
when the agent decides it applies, or when a person invokes it directly, so the
decision to bring material into context moves from a RAG pipeline or a static
file to the agent itself. But the decision stops at the skill boundary. Once
invoked, the entire \texttt{SKILL.md} lands in context, and a long one carries
the same situational content a long instruction file does. The agency costs the
whole file.

In January 2026 I kept some skill reference material as JSON and wrote the rows
of a \texttt{SKILL.md} as \texttt{jq} queries against those files. The agent ran
the query it needed and received the value. The file itself never entered
context. I call this context traversal without ingestion.

Most of my reference material is structured prose, though. A markdown file
carries a heading tree the way a JSON file carries keys, and what it lacked was a
command-line utility that would list that tree for an agent or return the prose
under one path through it. I wrote surf in March 2026 to be that utility. It
addresses a section by its heading and returns the section. Run against the
usage reference that ships with the surf skill:

\begin{Verbatim}
$ surf references/usage.md --list
- Surf Usage
  - Default: file map
  - Frontmatter only
  - Heading tree only
  - Targeted section
  - TeX
  - PDF
  - Level
  - Wikilink and link input
  - Batch scanning
  - Platform aggregation (batch section extraction)
  - Additional options

$ surf references/usage.md "Level"
## Level

`--level 1` is the top rank. On markdown, rank 1 is `#`. A file whose
first heading is `##` needs `--level 2` to list it. On TeX, rank 1 is
the shallowest command in that file. In an article that starts at
`\section`, `--level 1` is `\section`. On PDF, rank 1 is the top
outline rank.

`--list --level N` prints ranks 1 through N. A named heading with
`--level N` matches only a heading at rank N.
\end{Verbatim}

The agent has seen twelve lines of headings and one section from the middle of
the file, and it has not opened the file. By June I was building skill indexes
shaped by intent on top of this, and in September surf gained TeX sectioning
commands and PDF outlines so that the same skills could index academic papers.

\section{A Thin Index Shaped by Intent}

A skill built on surf has two parts. The \texttt{SKILL.md} is a short file: one
sentence stating what the skill does, followed by a table grouped by the intents
an agent might have on arrival. Each row pairs an intent with one command that
resolves to a section of a reference file by its heading.

\begin{Verbatim}
| Intent                      | Section extraction                            |
|-----------------------------|-----------------------------------------------|
| See what a file contains    | surf references/usage.md "Default: file map"  |
| Pull one section            | surf references/usage.md "Targeted section"   |
| Address a TeX paper         | surf references/usage.md "TeX"                |
\end{Verbatim}

Those three rows come from the surf skill's own index, and each resolves
against the \texttt{usage.md} that ships beside this paper.

The \texttt{references/} directory holds the corpus. I break reference files
into sections of twenty to eighty lines and write each so that an agent arriving
at it with no surrounding context can act from it alone.

The two parts have different jobs and change for different reasons. The index
is intensional. It defines, by intent, which sections matter. The corpus is
extensional. It holds the sections themselves. Datalog makes the same split
between an intensional database of rules and an extensional database of ground
facts \cite{ceri1989}, and the properties carry over. A row in the index holds a
query, never a copy of the content, and the query resolves against the corpus at
the moment the agent runs it. So the corpus can be rewritten or extended without
touching the index, and the index can be regrouped by intent without touching
the corpus. The only thing both sides depend on is the heading tree. The index
changes when the structure of the corpus changes, and at no other time.

A surf call with no heading returns that structure. \texttt{surf file.md}
prints the file's frontmatter and then its heading tree, and nothing from the
body. An agent's first move on an unfamiliar file is to look at the map, and its
second is to select a section. Neither move loads the file.

\section{Agentic Context Composition}

The question that separates this pattern from RAG is who decides what enters
the context window. In RAG a component outside the model decides, before the
model has seen anything. In a skill built on surf the model reads the intent
index, judges what its current task needs, and issues the pulls itself. Agent
inference, rather than a similarity score, composes the context.

Invocation loads the index and nothing else. This holds whether the agent chose
the skill or a person invoked it by name. When I tell an agent to surf my
holdings for papers on a topic, I have forced the skill to load, and I have
forced only the index. Which papers it lists, which sections it pulls, and in
what order remain the agent's decisions against the task I gave it.

The cost structure follows from this. In a long \texttt{SKILL.md}, presenting an
option to the agent costs the full text behind that option, whether or not the
agent takes it. In a thin index, presenting an option costs one row, and taking
it costs one section. A skill can therefore lay out every route it supports, and
an agent that needs one route pays for one section while an agent that needs
none pays for a table.

Models perform worse as the volume of irrelevant material in the context window
grows, even when the material they need is present \cite{liu2024,hong2025}.
Every token the model attends to that has no bearing on the task is a token of
inference spent on sorting rather than on working. A skill that hands the agent
only the sections it asked for spends that inference on the work.

\section{Skills You Can Check}
\label{sec:check}

Because the index is data, a script can verify the relationship between index
and corpus.

Every row is an executable command. After any edit to a skill, a script can run
every listed command and confirm that each resolves to a section. A row that
fails is a defect, and the most common cause is a renamed heading the index did
not follow. Renaming a heading is a schema change under this pattern, and the
check treats it as one.

The reverse check is more interesting. A script can list every section in the
corpus and count how many index rows reach each one. A section reached zero
times is a finding, and the count says which questions to ask. The index may
have fallen behind the corpus and be missing a row. The section may be content
nobody routes to, and belong trimmed or folded into a neighbor. The corpus may
have outgrown one index and need a second. Or the file may already be
structured for navigation by intent, so that an agent told to surf its heading
tree directly needs no row for it at all. Skills written as prose cannot be
audited this way, because there is nothing to count.

Two conventions make the checks reliable. The index quotes headings literally
from \texttt{surf -{}-list}, macros and punctuation included, so it never guesses
at an address. And structural divisions use real headings rather than bold text,
because surf sees the first and cannot see the second.

\section{Indexes over Indexes}

An index lists what a corpus contains. A procedure, a numbered list of steps,
says what to do with it. Most of what a skill carries is the first kind, and a
thin index hands the agent that inventory without ordering it. The order comes
from the task.

Indexes compose, and composing them is cheaper than enlarging one. A family of
twenty skills with forty sections each is eight hundred rows in a single index,
all of them loaded on every invocation to use one. Split into a top index that
routes by broad intent to member skills and a member index in each that routes
by narrower intent to its sections, the same destination costs twenty rows and
then forty, and the agent sees no others. A row in one skill's index may also
point at a section in another skill's corpus, so shared reference material lives
in one place and every skill that needs it routes there. The agent descends as
far as the task requires and at each level pays for one small table, never for
the corpora beneath it. A family grows by adding a member index and one row
above it.

\section{Markdown, TeX, and PDF}

surf addresses a markdown file by its ATX headings, a TeX file by its sectioning
commands and its \texttt{abstract} environment, and a PDF by its outline. The
agent's two moves are the same in each case: list the section titles, then
select the one to read.

\begin{Verbatim}
surf paper.tex --list --level 2
surf paper.tex "abstract"
surf paper.pdf "Experiments#Programming"
\end{Verbatim}

A markdown wiki adds a second kind of structure. In an Obsidian vault, a
section's body names other sections by link, \texttt{[[note\#Heading]]} or
\texttt{[[note\#Parent\#Child]]}, and surf accepts that link text as a target.
An agent reading one section follows a link with the same command it used to
arrive, so the vault is traversed one section at a time along the links its
authors wrote, and no note is loaded whole.

\begin{Verbatim}
surf "[[agents/memory#Episodic buffer]]"
\end{Verbatim}

A paper already has the structure the pattern needs, maintained by its authors.
One skill can therefore index a family of related papers at subsection
granularity, with rows that resolve to one experiment section of one paper. An
agent working the topic reads the sections its question touches across the
family and nothing else.

Shared structure across a corpus goes further. Most papers have an abstract, so
a loop over a directory of papers that runs \texttt{surf \$f "abstract"}
projects that one section out of most of the collection. If every skill in a
library keeps a \texttt{references/about.md} with the headings
\texttt{Overview} and \texttt{When to use}, two commands over the library return
a digest of every skill for a few hundred tokens each. Mine does, and the digest
is how I orient an agent to some eighty skills without loading one of them.
Frontmatter, read with \texttt{surf -f}, filters the collection to the files
that share a structure before the projection runs. A heading convention, seen
this way, is a schema declaration, and a corpus that follows one can be queried
like a table. In every format the address is the text that
\texttt{surf -{}-list} prints for the heading.

\section{This Paper}

The surf skill needs a TeX file and a PDF that it may ship, and this paper is
both. The commands below run against it as distributed.

\begin{Verbatim}
surf surf.tex --list
surf surf.pdf --list
surf surf.tex "abstract"
surf surf.pdf "Skills You Can Check"
\end{Verbatim}

The two listings return the same tree. The skill's build checks that claim
mechanically, in the way Section~\ref{sec:check} describes, each time this file
changes.

\begin{thebibliography}{9}

\bibitem{ceri1989}
S.~Ceri, G.~Gottlob, and L.~Tanca.
What you always wanted to know about Datalog (and never dared to ask).
\emph{IEEE Transactions on Knowledge and Data Engineering}, 1(1):146--166, 1989.

\bibitem{liu2024}
N.~F. Liu, K.~Lin, J.~Hewitt, A.~Paranjape, M.~Bevilacqua, F.~Petroni, and
P.~Liang.
Lost in the middle: How language models use long contexts.
\emph{Transactions of the Association for Computational Linguistics}, 12:157--173, 2024.

\bibitem{hong2025}
K.~Hong, A.~Troynikov, and J.~Huber.
Context rot: How increasing input tokens impacts LLM performance.
Chroma technical report, July 2025.

\end{thebibliography}

\end{document}
